\chapter{Why Parallel Computers?}
\label{ch:why}

\slides{2}

Every slide deck on parallel computing starts by saying that processors stopped getting faster.
That sounds odd: a laptop today is certainly faster than one from ten years ago. This chapter
explains what stopped, what did not, and why the industry's answer was to put \emph{many} processors on a chip
instead of one faster one.

\section{The free lunch is over}

For roughly thirty years, a programmer could write a program once and expect it to run faster on
every new generation of computer, with no changes at all. Two trends made this possible:

\begin{itemize}
  \item \textbf{More transistors.} A transistor is the tiny electronic switch that all digital
        circuits are built from. \term{Moore's law} is the observation that the number of
        transistors that fit on a chip doubles roughly every two years. This is still
        approximately true.
  \item \textbf{Higher clock frequency.} A processor does its work in steps called
        \emph{cycles}, paced by a clock. A $3$~GHz clock ticks three billion times per second.
        If you can raise the clock frequency, every instruction finishes sooner and every program speeds up.
\end{itemize}

The second trend is the one that ended. Clock frequencies climbed from a few MHz in the early
1980s to about $3$--$4$~GHz in the mid-2000s and then stopped, and they have stayed there. The
slide says it briefly: \emph{``CPU speeds have saturated''}. There are two main reasons.

\subsection{The power wall}

A chip turns almost all the electrical power it uses into heat. The \emph{dynamic} power of a
chip (the power spent switching transistors) is approximately
\begin{equation}
  P \;\approx\; \alpha\, C\, V^{2} f ,
  \label{eq:power}
\end{equation}
where $f$ is the clock frequency, $V$ the supply voltage, $C$ the total capacitance being
switched (roughly: how much circuitry there is), and $\alpha$ the fraction of the circuit that
switches in each cycle.

The catch is that $V$ cannot be chosen freely. A faster clock needs transistors to switch faster, which needs a
higher voltage. To a first approximation $V$ grows in proportion to $f$, so
\[
  P \;\propto\; f^{3}.
\]
Doubling the frequency would need about eight times the power, and that power has to leave the
chip as heat through an area the size of a postage stamp. Air cooling can remove on the order of
100~watts from a chip; beyond that you need exotic cooling that is too expensive for ordinary computers.
So frequency stopped rising.

\begin{workedex}{two slow cores beat one fast core}
Suppose one core runs at frequency $f$ and voltage $V$ and uses power $P_0$. Now run a core at
$0.8f$. Under the simplification $V \propto f$ it needs only $0.8V$, so its power is
\[
  P = 0.8 \times (0.8)^2 \, P_0 = 0.512\,P_0 .
\]
Two such cores together use $2\times 0.512\,P_0 \approx 1.02\,P_0$, essentially the same power as the
original single core, but if the work can be shared perfectly they deliver $2\times 0.8 = 1.6$
times the throughput. Making a chip with more, slower cores is therefore a far better use of a fixed power budget than
making a single core faster.
\end{workedex}

This is the central reason that parallel hardware became mandatory. The example hides an
important condition, though: ``\emph{if the work can be shared perfectly}''. The hardware
now offers parallelism; the software has to use it.

\subsection{The memory wall}

The second reason is that processors became much faster than main memory (DRAM). Two different
quantities matter here:

\begin{definition}{latency and bandwidth}
\textbf{Latency} is the time one request takes from start to finish, for example the time
between a core asking for a value and receiving it. \textbf{Bandwidth} is the amount of data
moved per unit time, for example gigabytes per second.
\end{definition}

\begin{intuition}
Think of a road. Latency is how long one car takes to drive from A to B. Bandwidth is how many cars
per hour the road can carry. A road can have low latency and low bandwidth (one fast lane) or the other way round
(a wide, slow motorway).
\end{intuition}

A main-memory access takes on the order of $100$~ns. At $3$~GHz that is hundreds of cycles during
which the core could have executed many instructions. Caches (your Chapter~1) hide much of this
latency by keeping recently used data close to the core. But caches do not help a program that
streams through a huge array once, and the memory \emph{bandwidth} of a chip is also finite and is
shared by all of its cores.

\begin{pitfall}
More cores do not automatically fix the memory wall. They bring more caches and more
opportunities to overlap computation with data movement, but they also compete for the same memory
bandwidth. This is one reason real machines have several memory channels and, as we will see in Chapter~2,
why large machines distribute their memory among processors (NUMA).
\end{pitfall}

\subsection{A third limit: instruction-level parallelism}

Inside a single core, hardware already runs several instructions at once when they do not depend on each other.
This is \term{instruction-level parallelism} (ILP): pipelining, several execution units, out-of-order
execution. It is invisible to the programmer. But finding independent instructions within one stream of ordinary code is
hard, and each extra bit of ILP costs a lot of hardware, so there are rapidly diminishing returns after a few instructions per cycle.
The remaining place to find parallelism is \emph{between} streams of instructions, which is
what multiple cores and multiple machines provide.

\section{Reasons to use parallel computing}

The slide lists four motivations. Each is worth stating in full.

\begin{enumerate}
  \item \textbf{Faster execution.} A computation that takes days or months on one core might
        take hours or seconds on many. Climate modelling and bioinformatics are the slide's examples:
        a weather forecast that arrives after the weather is useless.
  \item \textbf{Large amounts of data.} Some data sets do not fit in the memory of one machine,
        or take too long to scan with one processor. If a data set is 10~TB and one machine has 64~GB
        of RAM, the data must be spread over many machines and processed in parallel.
  \item \textbf{Natural parallelism.} Many problems are made of many similar, nearly independent parts.
        In a climate model the atmosphere is divided into a grid of cells, and the same update rule is applied to each
        cell. Writing this as a serial loop \emph{hides} parallelism that was in the problem all along.
  \item \textbf{Architecture trends.} The power and memory walls above: since a single core cannot be
        made much faster, the only way to use more transistors is to build more cores.
\end{enumerate}

There are two different goals hidden in these. One is to solve the \emph{same problem faster}; the
other is to solve a \emph{bigger problem in the same time}. Later slides give them precise names
(strong and weak scaling).

\section{Vocabulary used throughout}

Parallel computing uses a lot of words that sound alike. Fix the meanings now.

\begin{center}
\begin{tabularx}{\linewidth}{@{}l X@{}}
\toprule
\textbf{Term} & \textbf{Meaning in these notes}\\
\midrule
Core & One unit that executes one stream of instructions (possibly two when hyper-threading is used).\\
Processor / socket / chip & A package containing one or several cores. Slides often say ``processor''
  loosely to mean any processing element, so $P_1, P_2,\ldots$ may be cores or whole machines.\\
Node & One complete computer (``one box''): one or more sockets plus memory, disk and a network card.\\
Cluster & Many nodes connected by a network and used together. A supercomputer is in practice a very large cluster
  with a very fast network.\\
Process & A running program with its \emph{own} private address space.\\
Thread & An independent flow of execution \emph{inside} a process; all threads of a process share its memory.\\
Parallel & Several things really executing at the same instant (needs several cores).\\
Concurrent & Several things in progress, possibly interleaved on one core. All parallel programs are
  concurrent; not all concurrent programs are parallel.\\
\bottomrule
\end{tabularx}
\end{center}

\section{Not everything can be parallelised}

The slides remark that parallelism is ``more natural for certain kinds of problems''. The flip side is that
some work is inherently sequential: step~2 needs the result of step~1. If a fraction $s$ of a
program's running time is inherently serial, then no matter how many processors you use, the
speedup can never exceed $1/s$. A program that is $10\%$ serial can never run more than $10\times$
faster. The proper treatment (Amdahl's law) comes with the parallelisation-principles slides; keep this
number in mind.

\begin{keyidea}
Physics (power and memory limits) forced hardware designers to provide parallelism instead of
speed. Programmers now have to expose that parallelism in their programs. The rest of this part of the
course asks: what does parallel hardware look like, how do its parts talk to each other, and what does
that mean for performance?
\end{keyidea}

\section*{Exercises}
\addcontentsline{toc}{section}{Exercises}
\begin{enumerate}[label=\textbf{\thechapter.\arabic*},leftmargin=*]
  \item Using $P\propto f^3$, by what factor does the dynamic power of a core change if its clock frequency
        is raised by $25\%$?
  \item A core at frequency $f$ uses power $P_0$. Using the same simplified model, how much
        power do eight cores at $0.5f$ use in total, and what is their throughput relative to one core at $f$
        (assuming perfectly parallel work)?
  \item A program must read 12.8~GB from memory exactly once, and the memory system delivers
        25.6~GB/s. What is the shortest possible running time, regardless of how many cores are used or how
        low the latency is?
  \item A program is 5\% inherently serial. What is the largest possible speedup, and what speedup
        do 16 cores give? (Use $S(n)=1/(s+(1-s)/n)$.)
  \item Explain the difference between latency and bandwidth using an example other than a road.
\end{enumerate}